% Drop-in section fragment for iclr2026_cogen.tex.
% Required preamble: \input{math.tex}, natbib (loaded by iclr.sty), and hyperref/url.

\section{Preliminaries}
\label{sec:prelim}

We first specify the latent video backbone, the pixel-space action representation, and the
masked flow-matching objective shared by all task templates. We use $n$ for physical video time
and $k$ for an index in the discrete flow-matching schedule. A single noise level is sampled per
local forward pass and shared by every view and modality segment in the packed sequence. This
differs from UWM \citep{zhu2025unified}, which assigns independently sampled diffusion times to
its action and future-observation streams.

\subsection{Latent video backbone and camera conditioning}
\label{sec:prelim-backbone}

\paragraph{Video latent space.}
We instantiate the dense Wan2.2-TI2V-5B video diffusion transformer
\citep{wan2025,wan2025wan22}. Its causal 3D VAE encoder and decoder are denoted by $\gE$ and
$\gD$ and remain frozen. For admissible clips with $T=4q+1$ frames and spatial dimensions
divisible by $16$, the encoder maps
\begin{equation}
    \gE:\ \R^{3\times T\times H\times W}
    \longrightarrow
    \R^{C_l\times T_l\times H_l\times W_l},
    \qquad
    T_l=1+\frac{T-1}{4},\quad
    H_l=\frac{H}{16},\quad
    W_l=\frac{W}{16}.
    \label{eq:prelim-vae}
\end{equation}
For a general number of frames, the implementation uses
$T_l=1+\floor{(T-1)/4}$. The arm6 configuration uses $T=41$, $H=W=512$, and $C_l=48$, hence
$T_l\times H_l\times W_l=11\times32\times32$. We optimize all parameters of the denoising
transformer, including its camera-conditioning modules, while keeping the VAE and text encoder
fixed.

\paragraph{Absolute and relative cameras.}
Each example contains $V=2$ sampled views. Let
$\mK_{v,n}\in\R^{3\times3}$ and $\mE_{v,n}\in SE(3)$ be the intrinsics and absolute
camera-to-world extrinsics for view $v$ at frame $n$. They define the world-to-pixel projection
$\pi_{v,n}:\R^3\to\R^2$ by inverting $\mE_{v,n}$ before perspective projection. The arm6 cameras
are static within a trajectory, but we retain $n$ because the data interface supports per-frame
calibration.

The camera tensor supplied to the transformer is constructed from a different pose sequence.
Let $\mE_{0,1}$ denote the first frame of the first sampled (conditioning) view. We express every
camera in that coordinate system as
\begin{equation}
    \bar{\mE}_{v,n}=\mE_{0,1}^{-1}\mE_{v,n}=[\bar{\mR}_{v,n}\mid\bar{\vt}_{v,n}].
    \label{eq:prelim-relative-camera}
\end{equation}
For a pixel $\vx=(u,v)$, let
$\vd^{\mathrm{cam}}_{v,n,\vx}\propto\mK_{v,n}^{-1}(u+\tfrac12,v+\tfrac12,1)^{\top}$.
The corresponding world-frame ray has unit direction and Pl\"ucker moment
\begin{equation}
    \vd_{v,n,\vx}
    =\frac{\bar{\mR}_{v,n}\vd^{\mathrm{cam}}_{v,n,\vx}}
           {\|\bar{\mR}_{v,n}\vd^{\mathrm{cam}}_{v,n,\vx}\|_2},
    \qquad
    \vm_{v,n,\vx}=\bar{\vt}_{v,n}\times\vd_{v,n,\vx},
    \label{eq:prelim-plucker}
\end{equation}
which gives a six-channel ray map $(\vd,\vm)$ \citep{sitzmann2021lfn}. Following the camera
pathway of Action Images \citep{zhen2026action}, itself derived from camera-controlled video
generation \citep{bai2025recammaster}, we VAE-encode the direction and scaled-moment videos
separately and concatenate them:
\begin{equation}
    \tC^{(v)}
    =\operatorname{Concat}_{\mathrm{ch}}
      \left[\gE(\vd^{(v)}),\ \gE(\vm^{(v)}/3)\right]
    \in\R^{2C_l\times T_l\times H_l\times W_l}.
    \label{eq:prelim-camera-latent}
\end{equation}
The factor $1/3$ is part of the implementation and must be retained to reproduce the numerical
input scale. The camera latent for a view is repeated for every modality segment occupying
that view. At each packed latent-time position, every transformer block adaptively pools the
$2C_l\times H_l\times W_l$ camera descriptor to $48\times16\times16$, flattens it, projects it
to the hidden dimension, and adds it before self-attention. At module construction the encoder
projection is zero and the final attention projector is the identity, so the augmented block is
initially identical to the pretrained block. Arm6, however, is warm-started from the released
Action Images checkpoint; its optimization therefore begins from the camera weights stored in
that checkpoint rather than from the construction-time zero initialization.

The distinction between $\bar{\mE}_{v,n}$ and $\mE_{v,n}$ is essential: relative poses generate
the camera-conditioning tensor in \eqref{eq:prelim-camera-latent}, whereas absolute camera-to-world
extrinsics define $\pi_{v,n}$ for rendering and decoding action images.

\subsection{Actions as images}
\label{sec:prelim-actionimages}

An RLBench action at frame $n$ comprises an end-effector position $\vp_n\in\R^3$, a rotation
$\mR_n\in SO(3)$ obtained from the stored $xyz$ Euler angles, and a binary gripper state
$g_n\in\{0,1\}$. The renderer first constructs three world-space points. We write them using the
channel order implemented by the released code:
\begin{equation}
    \vq^{(R)}_n=\vp_n,
    \qquad
    \vq^{(G)}_n=\vp_n+\ell\mR_n\ve_x,
    \qquad
    \vq^{(B)}_n=\vp_n-\ell\mR_n\ve_z,
    \qquad \ell=0.1\ \mathrm{m},
    \label{eq:prelim-actionpoints}
\end{equation}
where $\ve_x=(1,0,0)^{\top}$ and $\ve_z=(0,0,1)^{\top}$. This is geometrically the same
three-point representation as Action Images \citep{zhen2026action}; the implementation permutes
the paper's semantic names for the green and blue axis points, and the decoder below follows the
implemented channel order.

For each view, the points are projected with the absolute calibration,
$\vu^{(c,v)}_n=\pi_{v,n}(\vq^{(c)}_n)$ for $c\in\{R,G,B\}$. Each projection is rasterized as an
isotropic Gaussian:
\begin{equation}
    \tilde A^{(v)}_{c,n,\vx}
    =\exp\!\left(-\frac{\|\vx-\vu^{(c,v)}_n\|_2^2}{2\sigma^2}\right),
    \qquad
    \sigma=0.05\min(H,W),
    \label{eq:prelim-actionimage}
\end{equation}
where $\vx$ ranges over pixel coordinates. Gripper openness is placed in the low-response
background of the blue channel, following Eq.~(5) of \citet{zhen2026action}:
\begin{equation}
    A^{(v)}_{B,n,\vx}=
    \begin{cases}
        \tilde A^{(v)}_{B,n,\vx},
        & \tilde A^{(v)}_{B,n,\vx}>\tfrac14,\\[2pt]
        \tfrac14\,\1_{g_n>1/2},
        & \tilde A^{(v)}_{B,n,\vx}\leq\tfrac14,
    \end{cases}
    \qquad
    A^{(v)}_{R,n}=\tilde A^{(v)}_{R,n},\quad
    A^{(v)}_{G,n}=\tilde A^{(v)}_{G,n}.
    \label{eq:prelim-openness}
\end{equation}
Thus $\tA^{(v)}\in[0,1]^{3\times T\times H\times W}$ is an ordinary three-channel video. It is
mapped to the frozen VAE range and encoded as
\begin{equation}
    \tZ^{(\mathrm{action},v)}=\gE(2\tA^{(v)}-1).
    \label{eq:prelim-action-latent}
\end{equation}
Rendering occurs inside the model forward pass: one calibrated 3D trajectory produces aligned
action videos in both views without additional per-view 2D action annotations.

\paragraph{Action decoding.}
The original Action Images decoder triangulates every point by selecting a heatmap centroid in a
principal view, sweeping depths along its camera ray, and scoring the reprojected candidates in
the remaining view. We retain that procedure for the red position point. For the two orientation
points, the arm6 evaluation and policy paths use the known axis length to avoid the degeneracy of
a depth sweep when an axis is foreshortened. Given a decoded position $\hat\vq^{(R)}_n$, we search
unit directions $\{\vd_j\}_{j=1}^{D}$ on the sphere and select
\begin{equation}
    \hat\vq^{(c)}_n
    =\hat\vq^{(R)}_n+\ell\vd_{j^*},
    \qquad
    j^*=\argmax_j\prod_{v=0}^{V-1}
    H^{(c,v)}_n\!\left(
        \pi_{v,n}(\hat\vq^{(R)}_n+\ell\vd_j)
    \right),
    \quad c\in\{G,B\}.
    \label{eq:prelim-sphere-decoder}
\end{equation}
The current solver uses $D=16384$ approximately uniform directions. Reported arm6 diagnostics
decode the position with $512$ depth samples over $[0.6,1.8]$ metres and use
\eqref{eq:prelim-sphere-decoder} for both axes. The paper's all-ray decoder remains available as a
separate baseline.

The gripper value is recovered from the blue-channel pedestal. To match the encoder's boundary
convention, define
$\Omega_n=\{(v,\vx):A^{(v)}_{B,n,\vx}\leq1/4\}$. The implementation uses
\begin{equation}
    \hat g_n
    =4\,\frac{1}{|\Omega_n|}
      \sum_{(v,\vx)\in\Omega_n}A^{(v)}_{B,n,\vx},
    \label{eq:prelim-openness-decoder}
\end{equation}
with a median fallback if $\Omega_n$ is empty. Finally,
\begin{equation}
    \hat\ve_x=\operatorname{norm}(\hat\vq^{(G)}_n-\hat\vq^{(R)}_n),\qquad
    \hat\ve_z=\operatorname{norm}(\hat\vq^{(R)}_n-\hat\vq^{(B)}_n),\qquad
    \hat\ve_y=\operatorname{norm}(\hat\ve_z\times\hat\ve_x),
    \label{eq:prelim-rotation-decoder}
\end{equation}
and $\hat\ve_z\leftarrow\hat\ve_x\times\hat\ve_y$ re-orthogonalizes the frame before forming
$\hat\mR_n=[\hat\ve_x\mid\hat\ve_y\mid\hat\ve_z]$.

\subsection{Packed sequences and masked flow matching}
\label{sec:prelim-objective}

Let $\Pi=(m_0,\ldots,m_{M-1})$ be the prompt-selected modality template in canonical order. The
latent sequence is packed view-major and modality-minor:
\begin{equation}
    \tZ
    =\big[\,\tZ^{(m_0,0)}\mid\cdots\mid\tZ^{(m_{M-1},0)}
      \mid\tZ^{(m_0,1)}\mid\cdots\mid\tZ^{(m_{M-1},1)}\,\big].
    \label{eq:prelim-packing}
\end{equation}
For the Action Images template this reduces to
$[\,V_0\mid A_0\mid V_1\mid A_1\,]$. Let $L$ be the packed latent-time length and
$\tZ_i=\tZ_{:,i,:,:}\in\R^{C_l\times H_l\times W_l}$. A sampled conditioning plan partitions
the latent-time positions into clean conditions $\sG\subseteq\{1,\ldots,L\}$ and prediction
targets $\sP=\{1,\ldots,L\}\setminus\sG$. A non-fully-given segment contributes its first latent
frame to $\sG$; a fully-given segment contributes all of its latent frames.

We use the optimal-transport linear path of flow matching \citep{lipman2023flow}. The training
schedule contains $K=1000$ discrete base noise levels
\begin{equation}
    \hat\sigma_k=1-\frac{k}{K},
    \qquad
    \sigma_k=\frac{s\hat\sigma_k}{1+(s-1)\hat\sigma_k},
    \qquad
    k\sim\operatorname{Unif}\{0,\ldots,K-1\},\quad s=5.
    \label{eq:prelim-shifted-schedule}
\end{equation}
Consequently, the grid runs from $\sigma_0=1$ to
$\sigma_{K-1}=0.005/1.004\approx0.00498$; it does not include zero. The scalar presented to the
Wan timestep embedding is $\tau_k=K\sigma_k$. For Gaussian noise
$\eps\sim\gN(\vzero,\mI)$ with the same shape as $\tZ$, we construct
\begin{equation}
    \tilde\tZ_k=(1-\sigma_k)\tZ+\sigma_k\eps,
    \qquad
    [\tilde\tZ_k]_i=\tZ_i\quad\text{for }i\in\sG,
    \qquad
    \tU=\eps-\tZ.
    \label{eq:prelim-noise}
\end{equation}
Thus predicted positions follow the linear data-to-noise path while conditions are restored to
their clean values.

The implementation additionally applies a scalar schedule weight. Define
\begin{equation}
    \bar w_k=\exp\!\left[-2\left(\frac{\tau_k-K/2}{K}\right)^2\right],
    \qquad
    w_k=K\,
    \frac{\bar w_k-\min_j\bar w_j}
         {\sum_{j=0}^{K-1}(\bar w_j-\min_r\bar w_r)}.
    \label{eq:prelim-timestep-weight}
\end{equation}
This gives mean weight one over the discrete schedule and places the largest weight near
$\sigma_k=1/2$. With prompt embedding $\vc$ and packed camera tensor $\tC$, the training loss is
\begin{equation}
    \gL(\vtheta)
    =\E_{\tZ,\sG,k,\eps}\!\left[
      \frac{w_k}{|\sP|C_lH_lW_l}
      \sum_{i\in\sP}
      \left\|
        \left[f_{\vtheta}(\tilde\tZ_k,\tau_k,\tC,\vc)\right]_i
        -\tU_i
      \right\|_2^2
    \right].
    \label{eq:prelim-loss}
\end{equation}
Condition positions receive no direct reconstruction loss, but their activations can receive
gradient through their influence on predicted positions. All modalities use the same output
projection, timestep distribution, and loss: there is no modality-specific prediction head or
loss coefficient. Task identity enters through the prompt and the contents and mask of the packed
sequence, not through a change to $f_{\vtheta}$.
